\chapter{VGA console in text mode}
\label{cap:vga}

\section{Overview}

The \code{src/vga.rs} driver is the simplified equivalent of Linux's
\code{vgacon}. It operates in the legacy VGA text mode, in which the screen is an
80-column by 25-row matrix stored starting from the physical address
\code{0xB8000}. Each cell occupies 16 bits.

\begin{table}[H]
\centering
\caption{Format of a VGA text cell.}
\label{tab:vgacell}
\begin{tabular}{@{}lll@{}}
\toprule
\textbf{Bits} & \textbf{Field} & \textbf{Value used} \\
\midrule
0--7  & ASCII character & Character code \\
8--11 & Foreground color & \code{0x7} (light gray) \\
12--14 & Background color & \code{0x0} (black) \\
15    & Blink & \code{0} (off) \\
\midrule
\multicolumn{3}{@{}l@{}}{\textbf{Full attribute}: \code{0x07}} \\
\bottomrule
\end{tabular}
\end{table}

The buffer has $80 \times 25 = 2{,}000$ cells, or 4,000 bytes. The cursor
position is controlled by the CRTC registers, accessed through ports
\code{0x3D4} (index) and \code{0x3D5} (data).

\section{Console state}

\begin{lstlisting}[caption={State maintained by the console.},label={lst:consolestate}]
struct Console {
    column: usize,
    row: usize,
    color: u8,
}
\end{lstlisting}

The state is protected by \code{static CONSOLE: SpinLock<Console>}. There is
no \emph{scrollback} buffer, nor memory separate from the screen: each
change is written directly to the \emph{framebuffer} with
\code{write\_volatile}.

\section{Operations}

\subsection{Initialization}

\code{vga::init()}:

\begin{enumerate}
  \item acquires the console lock;
  \item fills the 2,000 cells with a space and the \code{0x07} attribute;
  \item zeroes the column and row;
  \item configures the cursor as a bottom line, writing
        \code{0x0E} to register \code{0x0A} (cursor start) and
        \code{0x0F} to register \code{0x0B} (cursor end);
  \item updates the cursor position.
\end{enumerate}

\subsection{Byte output}

\begin{lstlisting}[caption={Handling of control and printable characters.},label={lst:writebyte}]
fn write_byte(&mut self, byte: u8) {
    match byte {
        b'\n' => self.newline(),
        b'\r' => { self.column = 0; self.update_cursor(); }
        b'\t' => {
            let next = (self.column + 8) & !7;
            if next >= BUFFER_WIDTH { self.newline(); }
            else { self.column = next; self.update_cursor(); }
        }
        0x08 => self.backspace(),
        0x20..=0x7E => self.write_char(byte as char),
        _ => self.write_char('?'),
    }
}
\end{lstlisting}

Printable ASCII characters are written directly. Any other
byte becomes \code{?}; this includes UTF-8 accented characters, which arrive as
individual bytes and are not decoded.

\subsection{Character output and line break}

\code{write\_char}:

\begin{itemize}
  \item if the column is at 80, it calls \code{newline()};
  \item computes the index \code{row * 80 + column};
  \item writes \code{(ch as u16) | (color << 8)} with \code{write\_volatile};
  \item increments the column;
  \item if the column has reached 80, it calls \code{newline()}; otherwise,
        it updates the cursor.
\end{itemize}

\code{newline()} zeroes the column and increments the row. If the current row is the
last one (24), it calls \code{scroll()}.

\subsection{Scrolling}

\code{scroll()} moves each cell from row $n$ to row $n-1$ and clears the
last row. The algorithm is $O(80 \times 24)$ per additional row and does not
use temporary memory:

\begin{lstlisting}[caption={Core of the scrolling.},label={lst:scroll}]
for row in 1..BUFFER_HEIGHT {
    for column in 0..BUFFER_WIDTH {
        let src = row * BUFFER_WIDTH + column;
        let dst = (row - 1) * BUFFER_WIDTH + column;
        let cell = read_volatile(VGA_BUFFER.add(src));
        write_volatile(VGA_BUFFER.add(dst), cell);
    }
}
/* limpa a ultima linha */
\end{lstlisting}

Because source and destination overlap and the copy is done from top to bottom,
there is no loss of cells that have not yet been copied.

\subsection{Backspace}

Backspace moves the column back; if it is at column 0, it moves to the end of the
previous row. It then writes a blank cell at the position and updates
the cursor. It does not erase beyond position (0,0).

\subsection{Hardware cursor}

The position is computed as \code{row * 80 + column} and split into two
bytes:

\begin{lstlisting}[caption={Cursor update via CRTC.},label={lst:cursor}]
outb(0x3D4, 0x0F);
outb(0x3D5, (position & 0xFF) as u8);
outb(0x3D4, 0x0E);
outb(0x3D5, (position >> 8) as u8);
\end{lstlisting}

Register \code{0x0E} receives the high byte and \code{0x0F} the low byte.

\section{Printing pipeline}

\begin{figure}[H]
\centering
\begin{tikzpicture}[
  box/.style={draw=rinuxblue, fill=rinuxlight, rounded corners=2pt,
              minimum width=2.8cm, minimum height=1.0cm,
              align=center, font=\scriptsize},
  arrow/.style={-{Latex[length=2mm]}, thick, rinuxblue},
  node distance=0.55cm
]
\node[box] (macro) {\code{println!}};
\node[box, right=of macro] (args) {\code{format\_args!}};
\node[box, right=of args] (print) {\code{vga::\_print}};
\node[box, right=of print] (writer) {\code{ConsoleWriter}};
\node[box, below=0.8cm of writer] (serial) {\code{serial::write\_str}};
\node[box, left=of serial] (vga) {\code{Console::write\_string}};
\node[box, left=of vga] (fb) {VGA \code{0xB8000}};
\draw[arrow] (macro) -- (args);
\draw[arrow] (args) -- (print);
\draw[arrow] (print) -- (writer);
\draw[arrow] (writer) -- (vga);
\draw[arrow] (writer) -- (serial);
\draw[arrow] (vga) -- (fb);
\end{tikzpicture}
\caption{Flow of a print call: VGA and serial receive the same text.}
\label{fig:printpipeline}
\end{figure}

\code{ConsoleWriter} implements \code{fmt::Write}. It acquires the lock,
writes to VGA, releases the lock, and only then writes to serial. This avoids
holding the lock while waiting for the UART.

\section{Public API}

\begin{table}[H]
\centering
\caption{Functions and constants of \code{vga.rs}.}
\label{tab:vgaapi}
\begin{tabularx}{\linewidth}{@{}lX@{}}
\toprule
\textbf{Item} & \textbf{Description} \\
\midrule
\code{BUFFER\_WIDTH = 80} & Screen columns. \\
\code{BUFFER\_HEIGHT = 25} & Screen rows. \\
\code{init()} & Clears the screen, configures the cursor shape and position. \\
\code{clear()} & Clears the screen (\code{clear} command). \\
\code{\_print(args)} & Entry point of the macros; writes to VGA and serial. \\
\code{Console::clear()} & Fills the buffer and repositions the cursor. \\
\code{Console::write\_byte()} & Interprets controls and characters. \\
\code{Console::write\_char()} & Writes a cell and advances the column. \\
\code{Console::write\_string()} & Iterates over the bytes of a \code{\&str}. \\
\code{Console::newline()} & Advances the row, scrolling if necessary. \\
\code{Console::backspace()} & Moves back and erases a cell. \\
\code{Console::scroll()} & Moves the content up by one row. \\
\code{Console::update\_cursor()} & Writes the position to the CRTC. \\
\bottomrule
\end{tabularx}
\end{table}

\section{Usage example}

\begin{lstlisting}[caption={Use of the macros in the shell.},label={lst:usovga}]
crate::print!("rinux> ");
println!("Comandos disponiveis:");
crate::print!("\x08");
\end{lstlisting}

The first command prints the prompt without a line break. The second prints
a complete line. The third sends the backspace character to erase the
last visible character.

\section{Limitations and observations}

\begin{itemize}
  \item \textbf{VGA text only}: there is no graphics framebuffer, VBE,
        GOP, or virtual console (VT). On UEFI without CSM, \code{0xB8000} may
        not exist.
  \item \textbf{ASCII only}: there is no UTF-8 decoding or
        built-in bitmap font; the font comes from VGA.
  \item \textbf{No dynamic attributes}: \code{color} exists in the state,
        but there is no public API to change it.
  \item \textbf{No synchronization with vertical retrace}: writes
        may cause occasional \emph{tearing}, irrelevant for text.
  \item \textbf{No memory protection}: \code{0xB8000} is accessed
        as a raw pointer; an invalid address would cause a fault.
  \item \textbf{Unchecked index}: \code{row} and \code{column}
        are kept within bounds by the functions, but there is no
        \emph{assert} of invariants.
  \item \textbf{Scroll costs 1,920 reads/writes}: acceptable for a
        terminal, but it could use \emph{hardware scrolling} or intermediate
        memory.
\end{itemize}
